% Part: incompleteness
% Chapter: theories-computability
% Section: computably-axiomatizable

\documentclass[../../../include/open-logic-section]{subfiles}

\begin{document}

\olfileid[hi]{inc}{tcp}{cax}

\olsection{\printtoken{S}{axiomatizable} सिद्धांत}

कोई सिद्धांत $\Th{T}$ \emph{!!{axiomatizable}} कहलाता है यदि उसके पास
अभिगृहीतों का कोई संगणनीय समुच्चय~$A$ हो। ($A$ को $\Th{T}$ के अभिगृहीतों
का समुच्चय कहने का अर्थ है
$T = \Setabs{!A}{A \Proves !A}$।) प्राकृतिक संख्याओं के प्रत्येक
``उचित'' अभिगृहीतीकरण में यह गुण होगा। विशेषतः, अभिगृहीतों के परिमित
समुच्चय वाला प्रत्येक सिद्धांत !!{axiomatizable} है।

\begin{lem}
मान लें $\Th{T}$,~!!{axiomatizable} है। तब $\Th{T}$,~!!{computably
enumerable} है।
\end{lem}

\begin{proof}
मान लें $A$,~$\Th{T}$ के अभिगृहीतों का कोई संगणनीय समुच्चय है। यह निर्धारित
करने के लिए कि $!A \in T$ है या नहीं, अभिगृहीतों से~$!A$ की किसी
!!a{derivation} की खोज करें।

कुछ भिन्न शब्दों में, $!A$,~$\Th{T}$ में तभी और केवल तभी है जब $A$ में
अभिगृहीतों $!B_1$, \dots, $!B_k$ की कोई परिमित सूची और प्रथम-कोटि
तर्कशास्त्र में $(!B_1 \land \dots \land !B_k) \lif !A$ की कोई
!!a{derivation} हो। किंतु हम पहले ही जानते हैं कि ``कोई \dots ऐसा है कि
\dots'' रूप की परिभाषा वाला प्रत्येक समुच्चय~!!{c.e.} है, बशर्ते दूसरा
``\dots'' संगणनीय हो।
\end{proof}

\end{document}
